\medskip
\paragraph{Formal Lean statement.}
\begin{verbatim}
private lemma residual_local_to_core_scaled
    (n t : ℕ) {c C₆ r₁ qSm qBig aFar cSm CSm cBig CBig cExp CExp : ℝ}
    (hn : 3 ≤ n) (hdim : 1 ≤ dim n) (ht_one : 1 ≤ t) (ht_pos : 0 < (t : ℝ))
    (hC₆_nonneg : 0 ≤ C₆) (hr₁_pos : 0 < r₁)
    (hbox : (dim n : ℝ) / (t : ℝ) ≤ Cn3Torus.delta ^ (2 : Nat))
    (hc_le_ann : c ≤ (1 / 32 : ℝ))
    (hc_le_sm : c ≤ cSm) (hc_le_big : c ≤ cBig) (hc_le_exp : c ≤ cExp)
    (ht_farSm : (t : ℝ) ≥ CSm * (n : ℝ) ^ (3 : Nat))
    (ht_farBig : (t : ℝ) ≥ CBig * (n : ℝ) ^ (3 : Nat))
    (ht_farExp : (t : ℝ) ≥ CExp * (n : ℝ) ^ (3 : Nat))
    (hannulus :
      ∀ (n : ℕ),
        1 ≤ dim n →
          ∀ {r delta : ℝ},
            0 < r →
              r ≤ r₁ →
                ∀ (t : ℕ),
                  1 ≤ t →
                    ∫ mu in (edgeBox n delta ∩ edgeEuclidBall n r) \ edgeCoreRegion n (t : ℝ),
                        ‖Cn3Torus.psi n mu‖ ^ (4 * t)
                      ≤ C₆ * Real.exp (-((dim n : ℝ) / 8)) * coreMass (dim n) (t : ℝ))
    (hfarShell :
      ∀ (n t : ℕ),
        2 ≤ n →
          1 ≤ t →
            Cn3Torus.texPrefactor n
              * ∫ mu in edgeEvenFarShell n r₁, ‖Cn3Torus.psi n mu‖ ^ (4 * t)
                ≤ qSm ^ (4 * t) + qBig ^ (4 * t)
                    + Real.exp (-(aFar * (t : ℝ) / ((n : ℝ) ^ (2 / 3 : ℝ)))))
    (hfarSm :
      ∀ n : ℕ, 3 ≤ n →
      ∀ t : ℕ, (t : ℝ) ≥ CSm * ↑n ^ (3 : Nat) →
        qSm ^ (4 * t) ≤ Real.exp (-(cSm * (n : ℝ) ^ (2 : Nat))) * gaussianScale n (t : ℝ))
    (hfarBig :
      ∀ n : ℕ, 3 ≤ n →
      ∀ t : ℕ, (t : ℝ) ≥ CBig * ↑n ^ (3 : Nat) →
        qBig ^ (4 * t) ≤ Real.exp (-(cBig * (n : ℝ) ^ (2 : Nat))) * gaussianScale n (t : ℝ))
    (hfarExp :
      ∀ n : ℕ, 3 ≤ n →
      ∀ t : ℕ, (t : ℝ) ≥ CExp * ↑n ^ (3 : Nat) →
        Real.exp (-(aFar * (t : ℝ) / ((n : ℝ) ^ (2 / 3 : ℝ))))
          ≤ Real.exp (-(cExp * (n : ℝ) ^ (2 : Nat))) * gaussianScale n (t : ℝ)) :
    |Cn3Torus.texPrefactor n * Cn3Torus.localIntegral n t - primaryCoreContribution n t|
      ≤ (C₆ + 3) * Real.exp (-(c * (n : ℝ) ^ (2 : Nat))) * gaussianScale n (t : ℝ)
\end{verbatim}
\begin{proof}[Proof sketch]
Combine the annulus and far-shell bounds and then transport the result from the local normalization to the return-probability scale. This is the last local-to-global assembly step before the final residual theorem.
\end{proof}
\paragraph{Role in the route.}
It depends directly on primaryCoreContribution. It is used directly by residual\_estimate\_quantitative.
\paragraph{Source location.}
\texttt{RequestProject/HadamardCn3LocalGapResidual.lean:253}
\medskip
